\documentclass[11pt]{article}

\usepackage[T1]{fontenc}
\usepackage[utf8]{inputenc}
\usepackage{lmodern}
\usepackage{amsmath,amssymb,mathtools}
\usepackage{booktabs}
\usepackage{enumitem}
\usepackage[margin=1in]{geometry}
\usepackage{microtype}
\usepackage[hidelinks]{hyperref}

\newcommand{\ind}{\mathbf{1}}
\newcommand{\given}{\,\middle|\,}
\DeclareMathOperator{\Close}{Close}
\DeclareMathOperator{\Open}{Open}
\DeclareMathOperator{\High}{High}
\DeclareMathOperator{\Low}{Low}
\DeclareMathOperator{\chol}{chol}
\DeclareMathOperator{\quantile}{quantile}
\DeclareMathOperator{\searchsorted}{searchsorted}

\title{BarrierLab: Mathematical Specification}
\author{}
\date{Measurement version \texttt{shared-outcome-cache-float64-v2}}

\begin{document}
\maketitle
\tableofcontents

\section{Scope}

This document specifies the numerical procedure implemented by the four-stage
BarrierLab pipeline. Caching and artifact references change where a tensor is
obtained, but they do not change any equation below.

\section{Notation and tensor axes}

Indices, sets, and sizes have one meaning throughout this document:
\begin{itemize}[nosep]
  \item $r\in\mathcal R$ identifies a price-history replicate; $r=0$ is the
        observed history and $r=1,\ldots,R$ identify null replicates;
  \item $i\in\{1,\ldots,N\}$ identify an observation date;
  \item $t\in\mathcal T$ identify a forward horizon;
  \item $\delta\in\mathcal D$ identify a signed return barrier;
  \item $k\in\{1,\ldots,K_{\mathrm{req}}\}$ identifies a requested bin
        position, while $K_r\leq K_{\mathrm{req}}$ is replicate $r$'s
        effective bin count; and
  \item $X_{r,i}$ denote a condition-feature value.
\end{itemize}

The default NASDAQ grid currently has $|\mathcal D|=41$,
$|\mathcal T|=30$, and requests $K_{\mathrm{req}}=10$ quantile bins.
Duplicate quantiles can make the effective count $K_r$ smaller than
$K_{\mathrm{req}}$. Batched tensors retain $K_{\mathrm{req}}$ positions and
mark positions above $K_r$ unsupported; compact observed artifacts retain
only $K_0$ bins. Stage~3 currently uses $R=1{,}000$ null replicates.

\begin{table}[ht]
\centering
\caption{Canonical tensor axes. An observed role batch has $|\mathcal R|=1$.}
\begin{tabular}{ll}
\toprule
Quantity & Axes \\
\midrule
OHLCV histories & $|\mathcal R|\times N\times 5$ \\
Forward downside/upside excursions & $|\mathcal R|\times N\times|\mathcal T|$ \\
Barrier touches & $|\mathcal R|\times N\times|\mathcal D|\times|\mathcal T|$ \\
Conditional probabilities & $|\mathcal R|\times|\mathcal D|\times K_{\mathrm{req}}\times|\mathcal T|$ \\
Baseline probabilities & $|\mathcal R|\times|\mathcal D|\times|\mathcal T|$ \\
Bin scores & $|\mathcal R|\times K_{\mathrm{req}}$ \\
\bottomrule
\end{tabular}
\end{table}

Synthetic histories are evaluated in memory-bounded path batches. The batch
axis is mathematically part of $\mathcal R$ and does not alter the statistic.

\section{Forward price excursions}

For every path, starting date, and horizon, define the greatest downside and
upside excursion reached by future intrabar prices:
\begin{align}
R^-_{r,i,t}
  &= \frac{\min_{1\leq j\leq t}\Low_{r,i+j}}
           {\Close_{r,i}}-1, \\
R^+_{r,i,t}
  &= \frac{\max_{1\leq j\leq t}\High_{r,i+j}}
           {\Close_{r,i}}-1.
\end{align}

Dates without $t$ future observations are non-finite and therefore
ineligible. Horizons are traversed incrementally: the running future minimum
and maximum for horizon $t$ extend the state calculated for shorter horizons.

\section{Signed-barrier touch matrix}

All barriers are broadcast over the excursion tensor. Define
\begin{equation}
Y_{r,i,t,\delta}
=
\begin{cases}
\ind\!\left[R^-_{r,i,t}\leq\delta\right], & \delta<0,\\[2pt]
\ind\!\left[R^+_{r,i,t}\geq\delta\right], & \delta\geq0.
\end{cases}
\end{equation}

For one replicate batch and horizon, $Y$ is a Boolean
$|\mathcal R|\times N\times|\mathcal D|$ matrix. It depends only on prices, horizon,
and the signed-barrier grid; it does not depend on a feature, node, or condition bin.
The implementation constructs it once and reuses it for the baseline and every
node. The zero barrier belongs to the measured probability surface but is
excluded from the paired-barrier score in Section~\ref{sec:score}.

\section{Feature bins}

For a continuous feature, the requested interior quantile edges are
\begin{equation}
e_{r,j}
=\quantile_{j/K_{\mathrm{req}}}\!\left(
  \{X_{r,i}:X_{r,i}\text{ is finite}\}\right),
\qquad j=1,\ldots,K_{\mathrm{req}}-1.
\end{equation}

Non-finite and duplicate edges are removed. Edges outside
$\min(X_r)<e_{r,j}\leq\max(X_r)$ are also removed. Heavy ties can therefore
produce $K_r<K_{\mathrm{req}}$ effective bins.

Bin assignment uses left insertion:
\begin{equation}
\kappa_{r,i}=\searchsorted(e_r,X_{r,i},\mathrm{left}).
\end{equation}
Consequently, an observation equal to an edge enters the lower bin. Stage~1
stores feature values, edges, and the resulting integer bin assignments.

For a synthetic OHLCV-derived feature, the feature and its quantile edges are
calculated independently for every synthetic path. External, calendar, and
workspace-extension conditions that cannot be reconstructed from synthetic
OHLCV use their observed feature values and observed edges for every null path.

\section{Conditional and unconditional probabilities}

Define price and condition eligibility by
\begin{align}
M^{\mathrm{price}}_{r,i,t}
  &=\ind\!\left[R^-_{r,i,t},R^+_{r,i,t}\text{ are finite}\right],\\
M^{\mathrm{condition}}_{r,i,t}
  &=M^{\mathrm{price}}_{r,i,t}
    \ind\!\left[X_{r,i}\text{ is finite}\right].
\end{align}

For condition bin $k$, horizon $t$, and barrier $\delta$, define
\begin{align}
N^{\mathrm{condition}}_{r,k,t}
  &=\sum_i M^{\mathrm{condition}}_{r,i,t}
    \ind[\kappa_{r,i}=k],\\
N^{\mathrm{hit}}_{r,\delta,k,t}
  &=\sum_i Y_{r,i,t,\delta}
    M^{\mathrm{condition}}_{r,i,t}\ind[\kappa_{r,i}=k],\\
\pi^{\mathrm{condition}}_{r,\delta,k,t}
  &=\frac{N^{\mathrm{hit}}_{r,\delta,k,t}}
          {N^{\mathrm{condition}}_{r,k,t}}.
\end{align}
The conditional probability is non-finite when
$N^{\mathrm{condition}}_{r,k,t}<30$.

The unconditional market baseline includes every price-eligible date and is
independent of feature warm-up:
\begin{align}
N^{\mathrm{baseline}}_{r,t}
  &=\sum_i M^{\mathrm{price}}_{r,i,t},\\
\pi^{\mathrm{baseline}}_{r,\delta,t}
  &=\frac{\sum_iY_{r,i,t,\delta}M^{\mathrm{price}}_{r,i,t}}
          {N^{\mathrm{baseline}}_{r,t}}.
\end{align}
The baseline is non-finite when $N^{\mathrm{baseline}}_{r,t}<30$. Every synthetic
path uses its own baseline; null histories are never compared with the
observed market's baseline.

\section{Stage 1: measurement}

Stage~1 measures the complete conditional surface
$\pi^{\mathrm{condition}}_{0,\delta,k,t}$, hit counts
$N^{\mathrm{hit}}_{0,\delta,k,t}$, observation counts
$N^{\mathrm{condition}}_{0,k,t}$, and the unconditional surface
$\pi^{\mathrm{baseline}}_{0,\delta,t}$. Observed forward excursions,
touches, and baseline probabilities are retained as float64 source-level cache
tensors. Feature values, quantile edges, and bin assignments are retained in
the node artifact.

\section{Stage 2: baseline-relative shift}

Stage~2 calculates the percentage-point deviation from the market baseline:
\begin{equation}
G_{r,\delta,k,t}
=100\left(\pi^{\mathrm{condition}}_{r,\delta,k,t}
          -\pi^{\mathrm{baseline}}_{r,\delta,t}\right).
\end{equation}

It stores only $G$ plus cryptographic references to the node and baseline
Stage~1 artifacts. Probabilities, counts, axes, conditions, and market history
are materialized from those referenced artifacts when required.

\section{Full-grid bin score}\label{sec:score}

Let $\mathcal M$ be the set of nonzero magnitudes $d$ for which the barrier grid
contains exactly one $+d$ barrier and one $-d$ barrier. For each pair, bin, and
horizon, define the signed-barrier asymmetry
\begin{equation}
A_{r,d,k,t}
=\left|G_{r,+d,k,t}-G_{r,-d,k,t}\right|.
\end{equation}
Each magnitude receives the linear weight
\begin{equation}
w_d=\frac{d}{\max_{m\in\mathcal M}m}.
\end{equation}

A contribution is usable only if
\begin{enumerate}[nosep]
  \item its positive and negative shifts are finite;
  \item its bin contains at least 30 eligible observations; and
  \item at least two bins are usable for that magnitude and horizon.
\end{enumerate}

Let $U_{r,d,k,t}$ be the corresponding Boolean usability indicator. The full
score for a condition bin is
\begin{equation}
Q_{r,k}
=\sum_{t\in\mathcal T}\sum_{d\in\mathcal M}
 U_{r,d,k,t}\,w_d\,A_{r,d,k,t}.
\end{equation}

An unsupported bin has score zero and a separate false-validity flag. A
supported bin whose contributions genuinely sum to zero remains valid.

\section{Synthetic null model}

For the observed history, define the four per-bar log components
\begin{align}
g^{\mathrm{open}}_i
  &=\log\!\left(\frac{\Open_i}{\Close_{i-1}}\right), &
g^{\mathrm{close}}_i
  &=\log\!\left(\frac{\Close_i}{\Open_i}\right),\\
g^{\mathrm{high}}_i
  &=\log\!\left(\frac{\High_i}{\max(\Open_i,\Close_i)}\right), &
g^{\mathrm{low}}_i
  &=\log\!\left(\frac{\Low_i}{\min(\Open_i,\Close_i)}\right).
\end{align}

After non-finite rows are removed, let $\mu$ and $\Sigma$ be the empirical
mean vector and covariance matrix. A small diagonal jitter is added before
Cholesky factorization:
\begin{align}
\epsilon
  &=10^{-12}\max\!\left(1,\max_j|\Sigma_{jj}|\right),\\
L_{\Sigma}&=\chol(\Sigma+\epsilon I_4).
\end{align}

Using deterministic seed \texttt{20260907}, each null innovation is
\begin{equation}
z_{r,i}\sim\mathcal N(0,I_4),
\qquad
g_{r,i}=z_{r,i}L_{\Sigma}^\mathsf{T}+\mu.
\end{equation}
The draws preserve contemporaneous covariance among the four OHLC components,
but they are independent across time.

Synthetic prices are reconstructed by
\begin{align}
\Close_{r,i}
  &=\Close_0\exp\!\left(
      \sum_{j\leq i}(g^{\mathrm{open}}_{r,j}+g^{\mathrm{close}}_{r,j})
    \right),\\
\Open_{r,i}
  &=\Close_{r,i-1}\exp(g^{\mathrm{open}}_{r,i}),\\
\High_{r,i}
  &=\max(\Open_{r,i},\Close_{r,i})\exp(g^{\mathrm{high}}_{r,i}),\\
\Low_{r,i}
  &=\min(\Open_{r,i},\Close_{r,i})\exp(g^{\mathrm{low}}_{r,i}).
\end{align}

Observed volume is broadcast across synthetic histories. The null does not
model temporal autocorrelation, volatility clustering, regime transitions,
liquidity, execution, or trading costs.

\section{Null scoring and hypothesis test}

Every synthetic history passes through the same excursion, touch, conditional
probability, baseline-shift, and bin-score equations used for the observation.
Within a path batch and horizon, the touch matrix is generated once; baseline
reduction and every node-specific bin reduction consume that same matrix.

For observed score $Q_{0,k}$ and $R$ null scores
$Q_{1,k},\ldots,Q_{R,k}$, the one-sided
Monte Carlo p-value is
\begin{equation}
\widehat p_k
=\frac{
  1+\sum_{r=1}^{R}
  \ind\!\left[Q_{r,k}\geq Q_{0,k}\right]
}{R+1}.
\end{equation}

With the current $R=1{,}000$, the denominator is $1{,}001$. A bin clears the
current raw threshold precisely when
\begin{equation}
\widehat p_k<0.05.
\end{equation}

Invalid null bins retain zero scores in the complete null ensemble. The output
separately records the number of null paths that produced a valid score. The
current decision rule applies no family-wise or false-discovery-rate correction
across bins.

\section{Stage 4 selection}

Stage~4 introduces no new statistical score. Let $C_k$ be the Stage~3
\texttt{cleared} indicator for condition bin $k$. The selected-bin set is exactly
\begin{equation}
\mathcal K_{\mathrm{sel}}
=\{k:C_k=\mathrm{true}\}=\{k:\widehat p_k<0.05\}.
\end{equation}
No top-$k$ truncation, family quota, effect-size gate, or representative-cell
rule is applied. Ordering by ascending $\widehat p_k$ and descending $Q_{0,k}$
is presentational only. For every $k\in\mathcal K_{\mathrm{sel}}$, Stage~4
renders the complete surface $G_{0,\delta,k,t}$ as a
signed-barrier-by-horizon heatmap.

\section{Cached mathematical identities}

The following quantities are reused exactly rather than approximated:
\begin{itemize}[nosep]
  \item observed forward-excursion tensors;
  \item observed signed-barrier touch matrices;
  \item observed unconditional baseline probabilities;
  \item observed feature values, quantile edges, and bin assignments;
  \item one synthetic touch matrix across the baseline and every node in a batch;
  \item common rolling feature primitives within a history or synthetic batch.
\end{itemize}

Observed caches are keyed by ordered market-history content and validated
against the measurement version, barrier grid, and horizon grid. Stage~2
references are protected by SHA-256 hashes. A changed history, grid, numerical
version, or referenced Stage~1 artifact invalidates reuse.

\end{document}
